\section{Introduction}

Agent systems are moving from one-shot task execution toward iterative improvement. Modern agents may revise prompts, update tool wrappers, adjust memory rules, change retrieval policies, modify planners, or promote new release candidates after receiving feedback. This shift creates an evidence problem for computational research infrastructure. A normal trace can record what happened during a single run, such as a model call, tool invocation, retrieved document, or output. A self-improvement event is different: it changes an artifact that may influence future runs. The accountability question therefore becomes temporal. A reviewer, downstream agent, or operator needs to know why a change was proposed, which prior version it derived from, what evidence supported it, how it was validated, why it was promoted or rejected, and how it can be rolled back if it later proves unsafe or invalid.

Existing work addresses important parts of this problem, but not the event-level evidence object itself. Self-Refine, Reflexion, and Voyager show that agents can improve behavior through feedback, reflection, or accumulated skills \cite{self_refine_2023,reflexion_2023,voyager_2023}. AgentDevel frames self-evolving agents as a release-engineering problem involving traces, candidate versions, and regression gates \cite{agentdevel_2026}. Observability systems such as OpenTelemetry GenAI and LangSmith provide execution traces, spans, runs, feedback records, and evaluation signals \cite{otel_genai_semconv,langsmith_observability_concepts}. Provenance and packaging frameworks such as W3C PROV, FAIR Digital Objects, and Workflow Run RO-Crate provide foundations for representing activities, entities, agents, digital objects, and machine-actionable research packages \cite{w3c_prov_dm,fdo_architecture_spec,workflow_run_rocrate}. These layers are valuable, but they do not define the minimum self-improvement-specific evidence required to audit one reusable agent change.

This paper proposes ASIEP, the Agent Self-Improvement Evidence Profile. ASIEP is a minimal FAIR-compatible profile for representing an agent self-improvement cycle as a verifiable, lifecycle-aware, reference-bound evidence object. It does not introduce a new self-improvement algorithm, agent framework, benchmark, or production governance system. Instead, it defines what must be present for a self-improvement event to be auditable: trigger evidence, lineage, candidate change evidence, validation results, gate decisions, integrity metadata, rollback references, and compliance metadata. ASIEP uses a reference-and-digest design so that raw prompts, traces, user inputs, model outputs, and gate artifacts do not need to be embedded directly in the core evidence object.

The paper makes three contributions. First, it defines the ASIEP profile, including field groups, lifecycle states, and invariants for agent self-improvement evidence. These invariants prevent common invalid states, such as promotion without a gate report, gating before a candidate is frozen, or evidence locking without digest-bound references. Second, it implements an agent-native local toolchain consisting of a validator, repair planner, resolver, importer, packager, and evaluator. The toolchain produces machine-readable validation errors, bounded repair plans, digest-checked bundle resolution, local trace imports, FDO-like and RO-Crate-like package outputs, PROV JSON-LD, and evaluation reports. Third, it evaluates the current implementation on a local fixture corpus containing valid records, invalid records, tampered bundles, missing artifacts, digest mismatches, path-escape cases, import fixtures, package requests, and adversarial evidence cases.

The scope is deliberately narrow. ASIEP is not a certified standard, not a production deployment, not a full OpenTelemetry or LangSmith integration, not a real FDO registry implementation, and not full RO-Crate certification. The evaluation is local and curated. Its purpose is to test whether a minimal evidence profile and local toolchain can support reproducible conformance checking, evidence closure, attack detection, and cross-standard mapping for agent self-improvement events.
